% Current mathematical front matter.  The historical review pages that follow
% are preserved as an audit record and are not silently rewritten.
\pdfbookmark[0]{Current review findings and proof status}{audit-corrective-addendum}
\thispagestyle{empty}
\begingroup
\raggedright
\setlength{\parindent}{0pt}
\setlength{\parskip}{0pt}
\setlength{\fboxsep}{5pt}
\newcommand{\ResultHead}[1]{\par\vspace{3.4pt}{\fontsize{9.15pt}{10.25pt}\selectfont\bfseries\color{AuditReviewed}#1}\par\vspace{1.2pt}}
\newcommand{\ResultSmall}{\fontsize{8.05pt}{9.23pt}\selectfont}

{\fontsize{15.8pt}{17.2pt}\selectfont\bfseries Independent Pro-mode reviews: the mathematical findings}\par
\vspace{1pt}
{\fontsize{7.75pt}{8.9pt}\selectfont\color{AuditProvenance}Reviewer B's exact reassessment is given here; Reviewer A's preserved independent corrective audit is on page~2, and the later Codex results and exact normalization certificates follow on pages~3--9.}
\par\vspace{1.2pt}
{\fontsize{7.25pt}{8.2pt}\selectfont\color{AuditProvenance}\textit{Editorial notation:}
$V=H^1(F;\mathbb Z)=\langle\gamma,u,w,\delta\rangle$,
$q=uw+6\gamma\delta\in\bigwedge^2V$, $\bar\delta$ is the class of
$\delta$ in $V_G$, and $\omega$ generates $H^2(B;\mathbb Z)$.}
\vspace{4pt}

\noindent\fcolorbox{AuditExact}{AuditPale}{%
\begin{minipage}{0.958\textwidth}
\ResultSmall
\textbf{Reviewer B's exact revised disposition: RECONSIDERED --- PROVISIONALLY
FAVORABLE WITH FOCUSED VERIFICATION REQUESTS.}
The earlier topology-based rejection is withdrawn.  Reviewer B found no
load-bearing defect in the intrinsic nearby-cycle/Leray and based van Kampen
route; it did not claim completed journal certification.
\end{minipage}}

\ResultHead{1. The correction is a distinction between two morphisms}
{\ResultSmall
Let $F$ be a nearby smooth torus fibre and $W$ the reduced cusp fibre.
The optional explicit collapse $c:F\to W$ would induce
\[
 c^*:H^q(W;\mathbb Z)\longrightarrow H^q(F;\mathbb Z),
\]
whereas nearby cycles supply the intrinsic stalk specialization
\[
 \operatorname{sp}^q:H^q(W;\mathbb Z)
 \longrightarrow H^q(F;\mathbb Z)^{T_0}.
\]
The earlier audit treated the uncompleted identification of $c^*$ with
specialization as a premise of both topology proofs.  Section~7.7 uses only
\(\operatorname{sp}^q\).  Reviewer B reconstructed Appendix~B's proof,
independent of $c$, that
\[
 \operatorname{sp}^q:H^q(W;\mathbb Z)
 \xrightarrow{\ \sim\ }H^q(F;\mathbb Z)^{T_0},\qquad 0\le q\le4,
 \qquad \operatorname{rk}=(1,2,4,2,1).
\]
Thus incomplete certification of the collapse map does not block the second
topology route.
}

\ResultHead{2. The same integer controls Leray and the based fundamental group}
{\ResultSmall
For the manuscript's twists
\[
 (\ell_0,\ell_1,\ell_2)=(0,1,-1),\qquad
 p=12\ell_0-4\ell_1-3\ell_2=-1.
\]
The reconstructed global stalk sections are
\[
 H^0(R^1f_*\mathbb Z)=\mathbb Z(12\gamma),\quad
 H^0(R^2f_*\mathbb Z)=\mathbb Z(2q),\quad
 H^0(R^3f_*\mathbb Z)=\mathbb Z(2\gamma uw),
\]
and Reviewer B obtained, up to a common orientation sign,
\[
 d_2(12\gamma)=\pm p\omega,\qquad
 d_2(2q)=\pm p\,\omega\bar\delta,\qquad
 d_2(2\gamma uw)=\pm p\,\omega[\gamma\delta].
\]
Because $p=-1$, all three rank-one maps are units, giving
\(H^1(X;\mathbb Z)=H^2(X;\mathbb Z)=H^3(X;\mathbb Z)=0\).  The based
overlap calculation gives
\[
 \pi_1(X)=\langle c,x,y\mid c\ {\rm central},\
 x^3=c^{\ell_1},\ y^4=c^{\ell_2},\ xy=c^{\ell_0}\rangle;
\]
its nonabelian elimination is cyclic of order $|p|$, hence trivial for
the selected twists.  Recognition then gives the integral homology and
smooth type of $S^6$.
}

\ResultHead{3. The algebraic-dimension computation is independent of topology}
{\ResultSmall
Reviewer B also reconstructed the period argument.  Off a countable set the
integral $(1,1)$-kernel is
\[
 \operatorname{NS}(F_b)=\mathbb Z\eta,\qquad
 \eta=uw+6\gamma\delta,\qquad \eta^2=12,
\]
and the corresponding Hermitian form has
\[
 \det G_\eta=\frac23\operatorname{Im}\tau
 \left(\operatorname{Im}\beta-
 \frac{6(\operatorname{Im}\mu)^2}{\operatorname{Im}\tau}\right)<0.
\]
The torus algebraic-reduction alternatives are thereby excluded; meromorphic
descent yields
\[
 \mathcal M(X)=f^*\mathbb C(t),\qquad a(X)=1,
\]
and identifies $f:X\to\mathbb P^1$ as the algebraic reduction, without
using Section~7 or the CDP comparison.
}

\vfill
{\fontsize{6.45pt}{7.38pt}\selectfont\color{AuditProvenance}
\textbf{Reviewer B's 26 August analytic follow-up.}  From a compact smooth
$X$ and nonconstant $f$, it independently proved flatness and, for every
scheme-theoretic surface fibre, $R^qf_*E=0$ for $q>2$ and
$(R^2f_*E)_b\otimes\mathbb C(b)\simeq H^2(X_b,E_b)$ by the finite-complex
universal-coefficient sequence.  From the asserted reduced Cartier cusp
fibre with normalization $\mathrm{dP}_6$ and $H^2(X;\mathbb Z)=0$, its
conductor--duality calculation then gives
$R^2f_*(T_X\otimes L)\ne0$ for every $L$.  The follow-up treats those three
candidate inputs as named hypotheses, so it independently confirms the
analytic arrows on page~3 without retrospectively certifying construction or
topology that it did not audit.  The seven-page report and eight exact local
certificates are indexed under
\nolinkurl{supporting_materials/peer_reviews/reviewer_B_candidate_validation_2026-08-26/}.}
\endgroup


\newpage
\thispagestyle{empty}
\begingroup
\raggedright
\setlength{\parindent}{0pt}
\setlength{\parskip}{0pt}
\setlength{\fboxsep}{5pt}
\newcommand{\ResultHead}[1]{\par\vspace{3.7pt}{\fontsize{9.2pt}{10.3pt}\selectfont\bfseries\color{AuditReviewed}#1}\par\vspace{1.3pt}}
\newcommand{\ResultSmall}{\fontsize{8.15pt}{9.42pt}\selectfont}

{\fontsize{15.2pt}{16.6pt}\selectfont\bfseries Reviewer A: corrective full-manuscript audit}
\par\vspace{2pt}
{\ResultSmall
Reviewer A withdrew the earlier packet-only rejection after auditing the
complete 108-page reconstruction.  Its revised disposition is
\textbf{MAJOR REVISION AND INDEPENDENT SPECIALIST REPLICATION}: the algebra
and two non-collapse arguments are substantially stronger than the first
report allowed, but the geometric and analytic bridges were not promoted
beyond the evidence seen at its 25 August cutoff.
}

\ResultHead{Exact algebra and the topology implication}
{\ResultSmall
Two independent harnesses reproduced 28 legacy and 14 supplemental checks:
the monodromy invariant/coinvariant lattices, Smith forms, period laws,
$P_E$, $\eta^2$, the Gram determinant, and the Appendix~B rank and
saturation matrices.  The topology depends on
\[
p=12\ell_0-4\ell_1-3\ell_2=-1
\quad\text{for}\quad(\ell_0,\ell_1,\ell_2)=(0,1,-1).
\]
Given the actual based meridians and drift terms, van Kampen yields
$\pi_1(X)\cong\mathbb Z/|p|$.  Independently, if the integral cusp and
finite-fibre specializations are the geometric stalk maps, then
\[
d_2(12\gamma)=\pm p\omega,\quad
d_2(2q)=\pm p\omega\bar\delta,\quad
d_2(2\gamma uw)=\pm p\omega[\gamma\delta].
\]
Thus $p=-1$ makes all three transgressions units and implies the integral
homology and smooth type of $S^6$.  Reviewer A's correction is precise:
the explicit collapse is not theorem-critical for this Leray route, although
the identified meridian, specialization, and transgression bridges still
required specialist replication at its cutoff.
}

\ResultHead{The algebraic-dimension proof survives independently}
{\ResultSmall
For an integral class $E$, Section~9.1 derives
$E\in\operatorname{NS}(F_b)\Longleftrightarrow P_E=0$ and proves the linear
independence of $1,\tau,\mu,\beta,6\mu^2-\tau\beta$.  Hence, off a countable
set,
\[
\operatorname{NS}(F_b)=\mathbb Z\eta,\qquad
\eta=uw+6\gamma\delta,\qquad \eta^2=12,
\]
while
\[
\det G_\eta=\frac23\operatorname{Im}\tau
\left(\operatorname{Im}\beta-
\frac{6(\operatorname{Im}\mu)^2}{\operatorname{Im}\tau}\right)<0.
\]
There is therefore neither an ample positive class nor a nonzero isotropic
fibre class, excluding algebraic dimensions two and one for the very general
torus fibre.  Meromorphic descent gives
$\mathcal M(X)=f^*\mathbb C(t)$ and $a(X)=1$.  Reviewer A classified this
argument as complete conditional only on the period-family construction and
standard complex-torus structure theory; it uses neither Section~7 nor CDP.
}

\ResultHead{What the Section 10 calculation does and does not prove}
{\ResultSmall
On the nonnormal cusp fibre, the ambient differentials in the local models
$xy=0$ and $xyz=0$ have coefficient vectors $(y,x,0)$ and $(yz,xz,xy)$.
They vanish on the conductor and descend to a nonzero section of
$\Omega_X^1|_W\otimes A$ whose intrinsic image is conductor-supported
torsion.  Consequently, vanishing after passage to the torsion-free
$\widetilde\Omega_W^1$ does not imply vanishing of the ambient cotangent
group---the exact inference used in the printed CDP comparison.  Reviewer A
found this local mechanism coherent, but kept separate (i) a defect in that
printed inference, (ii) falsity of the published theorem, and (iii) validity
of the candidate.  At its cutoff, the global conclusion still required the
stated Gorenstein duality, top-degree base change, Picard input, and complete
construction of $X$.
}

\ResultHead{Confirmed local corrections}
{\ResultSmall
The audit confirms four exact repairs: Proposition~4.7(iv) needs the
closed-graph/section argument rather than ``proper over an open disc'';
Remark~7.28 is existential in $\ell_0$, not a biconditional for a fixed
triple; Lemma~A.3's displayed Mayer--Vietoris sequence is cohomology and
must be followed by universal coefficients; and the cusp deck group is
$\bar\Lambda=\Lambda/\Lambda_{\rm tor}$.  The modular factor
$\mathsf e=F^{-1}e$ on source page~71 was already present in the corrected
source and is not an additional current defect.
}

\vfill
{\fontsize{6.7pt}{7.7pt}\selectfont\color{AuditProvenance}
Reviewer A's full corrective report, exact-check receipts, proof-obligation
ledger, and CDP comparison are indexed under
\nolinkurl{supporting_materials/peer_reviews/reviewer_A_corrective_2026-08-25/}.
The page records that review's own evidentiary boundary; it is not a summary
of the later Codex proofs on page~3.}
\endgroup

\newpage
\thispagestyle{empty}
\begingroup
\raggedright
\setlength{\parindent}{0pt}
\setlength{\parskip}{0pt}
\setlength{\fboxsep}{5pt}
\newcommand{\ResultHead}[1]{\par\vspace{2.1pt}{\fontsize{8.25pt}{9.05pt}\selectfont\bfseries\color{AuditReviewed}#1}\par\vspace{0.7pt}}
\newcommand{\ResultSmall}{\fontsize{7.12pt}{8.12pt}\selectfont}

{\fontsize{15.2pt}{16.6pt}\selectfont\bfseries Subsequent Codex validate/repair/disprove results}
\par\vspace{1pt}
{\fontsize{7.1pt}{8.1pt}\selectfont\color{AuditProvenance}
These are later workbench proof classifications, presented separately so
they are not attributed retrospectively to Reviewer A or Reviewer B.  Each
item is \textbf{VERIFIED} in the cited proof file as of 30 August 2026.}

\ResultHead{Why the result matters}
{\ResultSmall
An integrable complex structure on the smooth six-sphere would settle the
classical existence problem.  This candidate is explicit enough to audit in
coordinates and would also give a compact complex threefold with
\(b_2=0\), \(a(X)=1\), and \(c_3(X)=2\), contrary to CDP20 Theorem~2.2 and
Corollary~2.3 as stated.
}

\ResultHead{Global analytic threefold and all punctured-overlap maps --- VERIFIED}
{\ResultSmall
The reconstructed periods, finite logarithmic quotients, toric cusp quotient,
and punctured-disc identifications give a connected compact Hausdorff complex
threefold and a proper map \(f:X\to\mathbb P^1\).  At the cusp the covering is
\[
 (z,\zeta_1,\zeta_2)\longmapsto
 (e^{2\pi i\zeta_1},e^{2\pi i\zeta_2},e^{2\pi i s}),
\]
with deck kernel \(\Lambda_{\rm tor}\times\langle g_0\rangle\) and residual
deck group \(\bar\Lambda\); the special fibres are \(W,3S_1,4S_2\).
The finite signs are one simultaneous orientation normalization,
\[
 g_1'(z_1)=e^{-2\pi i/3},\qquad g_2'(z_2)=-i,
\]
and the exact modular deck certificate is
\(\varphi(g_1)=S_1,\varphi(g_2)=S,\varphi(g_0)=P\), with
\(S_1SP=1\).  Quadratic ramification alone gives only
\((g_2'(z_2))^2=-1\).
}

\ResultHead{Intrinsic specialization and the Leray differentials --- VERIFIED}
{\ResultSmall
Integral nearby cycles give the actual cusp and finite-fibre stalk maps,
including primitivity and saturation.  For \(K=Rf_*\mathbb Z_X\), the
canonical truncation triangle
\[
 R^0f_*\mathbb Z_X\longrightarrow\tau_{\le1}K
 \longrightarrow R^1f_*\mathbb Z_X[-1]
 \xrightarrow{\kappa}R^0f_*\mathbb Z_X[1],
\]
identifies the annular lift-discrepancy cocycle with Leray \(d_2\).  Writing
\(q:=uw+6\gamma\delta\in H^2(F;\mathbb Z)\) and
\(p=12\ell_0-4\ell_1-3\ell_2=-1\), the three maps are
\[
 d_2(12\gamma)=\omega,\qquad
 d_2(2q)=\omega[\delta],\qquad
 d_2(2\gamma uw)=\omega[\gamma\delta],
\]
so \(H^1=H^2=H^3=0\) and \(H_2=0\).
}

\ResultHead{Based van Kampen and sphere recognition --- VERIFIED}
{\ResultSmall
The actual based overlap maps, with fixed whiskers and meridians, give
\[
 \pi_1(X)=\langle c,x,y\mid c\ {\rm central},\ xy=1,\ x^3=c,\ y^4=c^{-1}\rangle=1.
\]
The cusp meridian bounds the toric disc and there is no omitted fibre drift.
Universal coefficients, Poincar\'e duality, Hurewicz, Whitehead, and
\(\Theta_6=0\) then give
\[
 X\cong_{\rm diff}S^6,
 \qquad \langle c_3(T_X),[X]\rangle=\chi(X)=2.
\]
}

\ResultHead{Algebraic reduction --- VERIFIED}
{\ResultSmall
The Appell--Humbert calculation gives, for very general \(b\),
\(\operatorname{NS}(F_b)=\mathbb Z(uw+6\gamma\delta)\) with indefinite Gram
form.  Hence \(a(F_b)=0\), and meromorphic descent gives
\[
 \mathcal M(X)=f^*\mathbb C(t),\qquad a(X)=1,
\]
with \(f\) the algebraic reduction, independently of topology and CDP.
}

\ResultHead{Proposition 10.8: uniform non-torsion --- VERIFIED}
{\ResultSmall
Exact direct images and base change at every scheme-theoretic fibre give
\[
 f_*\mathcal O_X=\mathcal O_B,\qquad
 R^1f_*\mathcal O_X\simeq\mathcal O_B\oplus\mathcal O_B(-1),\qquad
 R^2f_*\mathcal O_X\simeq\mathcal O_B(-1),\qquad
 \operatorname{Pic}(X)\simeq\mathbb C.
\]
Moreover \(\operatorname{ord}(\mathcal O_X(S_j)|_{S_j})=m_j\),
\((m_1,m_2)=(3,4)\).  For
\(\lambda=a\gamma+bu+cw+d\delta\), restriction to the global Picard line is
\[
 c\tau+d\mu-b=0.
\]
Independence of \(1,\tau,\mu\) leaves \(\mathbb R\) and a countable exceptional
set, so one line bundle is non-torsion on every full and reduced fibre.  The
independent equivariant-collapse/Mayer--Vietoris comparison is also complete:
its primitive integral maps commute with all three Leray transgressions.
}

\ResultHead{The Hodge calculation left open by the manuscript --- VERIFIED}
{\ResultSmall
For \(D=W+S_1+S_2\), cyclic invariant lattices, the toric cusp frame, and
normalized residues give
\[
R^0f_*\Omega_X^1=\mathcal O(-2),\quad
R^1f_*\Omega_X^1=\mathcal O\oplus\mathcal O(-1)^2,\quad
R^2f_*\Omega_X^1=\mathcal O\oplus\mathcal O(-1).
\]
The same local module calculation corrects the manuscript's torsion claim:
$\Omega^1_{X/B}$ is torsion-free but non-locally-free at the cusp, and its
torsion occurs only at the two multiple fibres.
The logarithmic Kodaira--Spencer determinant has divisor \(p_1+p_2\); at the
order-four lift \(\dot\tau=0\) but
\(6\dot\mu^2=\det(\dot Z)\ne0\), so
\[
h^{1,2}=h^{2,1}=1,\qquad h^{1,1}=h^{2,2}=2,\qquad E_2=E_\infty.
\]
\href{https://www.rivmat.unipr.it/vols/2019-10-1/02-mchugh.html}%
{McHugh's threefold identities} then give the total Bott--Chern and Aeppli
dimensions
\[
(h^k_{BC})=(1,0,2,2,2,2,1),\qquad
(h^k_A)=(1,2,2,2,2,0,1),
\]
so their non-$\partial\bar\partial$ degrees are
$(0,2,4,4,4,2,0)$.  Global representatives with exact marked-fibre
period-coordinate classes give nine positive-degree Bott--Chern generators
together with the unit \(\mathbf1=[1]_{BC}\).
Writing $N$ for their span, the complete algebra is
$H^{*,*}_{BC}(X)=\mathbb C\mathbf1\oplus N$ with $N^2=0$.  A preferred
literal global coefficient atlas through the fillings and cusp is not claimed.
}

\ResultHead{The cusp obstruction: Pic$^0$ and all-Pic restriction routes --- VERIFIED}
{\ResultSmall
The cusp form \(dx_1\wedge dx_2\wedge dt/(x_1x_2)\) trivializes \(\omega_X\)
near \(W\).  Let \(\nu:\widetilde W\simeq\mathrm{dP}_6\to W\) be the
normalization.  For \(L\in\operatorname{Pic}^0(X)\), naturality and injectivity
of \(c_1\) on \(\operatorname{Pic}(\mathrm{dP}_6)\) give
\[
 \nu^*((L^\vee\otimes\omega_X)|_W)\simeq\mathcal O_{\mathrm{dP}_6}.
\]
Conductor sections, Ramis--Ruget duality, and top-degree base change imply
\(R^2f_*(T_X\otimes L)\ne0\).  For arbitrary \(L\), the sole extra obstruction
is \(r=\nu^*i^*:\operatorname{Pic}(X)\to\operatorname{Pic}(\mathrm{dP}_6)\).
For this \(q\), the Leray edge image of \(c_1(L)\) is a permanent cycle, but
\(d_2^{0,2}(2q)=\omega[\delta]\) is an isomorphism; hence the edge section
vanishes.  Proper base change gives \(c_1(L|_W)=0\), so injectivity of
\(c_1\) on \(\mathrm{dP}_6\) gives \(r(L)=0\).  Independently of global
\(d_2\), the order-four filling writes \(c_1(L|_F)=2nq\).  The
pair-of-pants boundary defects are
\[
 \bar d_1=4n[\delta],\qquad \bar d_2=-3n[\delta],\qquad \bar d_0=0;
\]
the finite values are logarithmic-cover norm maps and the cusp value is a
literal bounding solid torus.  Their based sum is zero, so \(n=0\); local
specialization again gives \(r(L)=0\).  Thus the relative CDP sufficient
condition fails for every \(L\); separately, \(a(X)=1,c_3(X)=2\) give a
counterexample to CDP20
Theorem~2.2 as stated.  A third proof now avoids both global integral Leray
and local integral specialization.  Put \(M=L|_{N_0}\).  For
\(c_1(M|_{F_t})=kq\), its holomorphic
factor defects are
\[
 (\bar d_1,\bar d_2,\bar d_0)=(2k,-\tfrac32k,0)[\delta],
 \qquad \bar d_1+\bar d_2+\bar d_0=0,
\]
so \(k=0\).  The cusp equality is an argument-principle calculation in a
Stein chart.  Thus the Appell--Humbert form
\(E_M:=c_1(M|_{F_t})=c_1(L|_{F_t})\) is zero.  With
\[
 B_0:\overline\Lambda\xrightarrow{\ \sim\ }\Lambda_{\rm tor},\qquad
 [B_0]_{(\overline{\widehat\gamma},\overline{\widehat u})}
 ^{(\widehat w,\widehat\delta)}
 =\begin{pmatrix}0&1\\-1&0\end{pmatrix},
\]
the normalized central degrees determine
\[
 B=\begin{pmatrix}d_2+d_3&d_3\\d_3&d_1+d_3\end{pmatrix},
 \qquad E_M((\lambda,0),(0,m))=-B(B_0\lambda,m)
 \quad(\lambda\in\overline\Lambda,\ m\in\Lambda_{\rm tor});
\]
the antecedent \(E_M=0\) and unimodularity of \(B_0\) force \(B=0\), all
six central degrees vanish, and hence \(r(L)=0\).
}

\vfill
{\fontsize{6.35pt}{7.2pt}\selectfont\color{AuditProvenance}
Complete proofs are in \nolinkurl{supporting_materials/workbench/research/},
especially \texttt{candidate\_geometry.tex}, \texttt{integral\_leray.tex},
\texttt{topology\_reduction.tex}, \texttt{algebraic\_dimension.tex},
\texttt{analytic\_picard.tex}, \texttt{analytic\_picard\_completion.tex},
\texttt{hodge\_frolicher.tex}, and
\texttt{cdp\_program.tex}.  External specialist replication remains open.\par
The earlier review and page-by-page annotations follow as a preserved
historical record.  Where an older status label conflicts with this front
matter, compile the root entry point \texttt{workbench.tex}.  Both topology
certificates and the Hodge calculation are complete there; independent
specialist replication remains open.}
\endgroup

\newpage
\thispagestyle{empty}
\begingroup
\raggedright
\setlength{\parindent}{0pt}
\setlength{\parskip}{0pt}
\setlength{\fboxsep}{5pt}
\newcommand{\ResultHead}[1]{\par\vspace{2.2pt}{\fontsize{8.25pt}{9.05pt}\selectfont\bfseries\color{AuditReviewed}#1}\par\vspace{0.7pt}}
\newcommand{\ResultSmall}{\fontsize{6.85pt}{7.82pt}\selectfont}

{\fontsize{14.6pt}{16pt}\selectfont\bfseries Exact sign, Hodge, and Bott--Chern certificates}
\par\vspace{1pt}
{\fontsize{6.8pt}{7.75pt}\selectfont\color{AuditProvenance}
This page records the exact conventions behind the compact conclusions on
page~3.  It is current workbench mathematics, not a retrospective attribution
to either independent reviewer.}

\ResultHead{The period signs are simultaneous}
{\ResultSmall
Let
\[
 S_1(w)=\frac{w-1}{w},\qquad S(w)=-\frac1w,\qquad P(w)=w-1.
\]
With \(g_0=(g_1g_2)^{-1}\), the normalized choices are
\[
 g_1'(z_1)=e^{-2\pi i/3},\qquad g_2'(z_2)=-i,
 \qquad g_0\text{ parabolic},
\]
and
\[
 \varphi(g_1)=S_1,\qquad\varphi(g_2)=S,\qquad
 \varphi(g_0)=P,\qquad S_1SP=1.
\]
The two derivative signs are one orientation choice: once the \(g_1\)
orientation and primitive parabolic cusp class are fixed, the mixed choice is
not parabolic.  The local quadratic calculation at \(z_2\) supplies only the
square of the derivative, not the sign.
}

\ResultHead{The invariant class and the two fibre orientations}
{\ResultSmall
In the marked basis \(V=\mathbb Z\langle\gamma,u,w,\delta\rangle\), put
\[
 y=\operatorname{Im}\tau,\qquad m=\operatorname{Im}\mu,\qquad
 r=\operatorname{Im}\beta,\qquad
 \Delta=r-\frac{6m^2}{y}<0,
\]
\[
 \Gamma=\frac{i}{2\Delta}
 \left(d\bar\zeta_2-\frac{m}{y}d\bar\zeta_1\right),\qquad
 U=\frac{i}{2y\Delta}
 \left(r\,d\bar\zeta_1-6m\,d\bar\zeta_2\right).
\]
Then
\[
 \eta=u\wedge w+6\gamma\wedge\delta,
 \qquad \eta^2=12\gamma\wedge u\wedge w\wedge\delta,
\]
and its exact marked-fibre Dolbeault representative is
\[
 Q_F=-d\zeta_1\wedge U-6d\zeta_2\wedge\Gamma.
\]
The determinant of the four marked lattice columns in real coordinates is
\(y\Delta<0\), so the marked and complex orientations are opposite.  Hence
\[
 \int_{(F_b,{\rm or}_{\rm mark})}Q_F^2=12,
 \qquad
 \int_{(F_b,{\rm or}_{\mathbb C})}Q_F^2=-12,
 \qquad
 \int_{(F_b,{\rm or}_{\mathbb C})}
 Q_F\wedge\Gamma\wedge d\zeta_2=-1.
\]
The degree-one class \(\mathfrak g\) is distinguished, whereas
\(\mathfrak q\) is a chosen coherent-Leray lift of the nonzero section
\(\boldsymbol\eta\in H^0(B,R^1f_*\Omega_X^1)\), and
\[
 i_b^*\mathfrak g=[\Gamma]_{\bar\partial},\qquad
 i_b^*\mathfrak q=[Q_F]_{\bar\partial}.
\]
The lift \(\mathfrak q\) is unique only modulo \(\operatorname{im}d_1\),
which restricts trivially to every smooth fibre.  Thus a smooth global
representative \(Q\) need not restrict coefficient-by-coefficient to \(Q_F\).
}

\ResultHead{The Kodaira--Spencer, residue, and cusp maps}
{\ResultSmall
For \(t=s^m\), \((m,a)=(3,2)\) or \((4,1)\), and
\(D_{\rm KS}=(d_{rs})\), put
\[
 A=\frac{s}{m}d_{00},\quad B=\frac{s^{1-a}}m d_{01},\quad
 C=\frac{s^{1-a}}m d_{10},\quad
 D_0=\frac{s^{m+1-2a}}m d_{11}.
\]
In the actual invariant bases the exact coarse matrices are
\[
 [\kappa_0]=\begin{pmatrix}A&tB\\C&D_0\end{pmatrix},
 \qquad [\kappa_1]=(-C,A,-D_0,B).
\]
At the cusp, in its separately fixed logarithmic frames,
\[
 [\kappa_0]\equiv\begin{pmatrix}0&-1\\1&0\end{pmatrix}\pmod t,
 \qquad [\kappa_1]\equiv(-1,0,0,-1)\pmod t.
\]
At a finite fibre the typed residue calculation is
\[
 \operatorname{res}_{1,j}(xb_0+ys^ab_1)=m_jx[b_0],\qquad
 \operatorname{im}(\kappa_0\bmod t)
 =\ker\operatorname{res}_{1,j}=\mathbb C\,s^ab_1.
\]
At the cusp take
\(q_1=(-1,0,1),q_2=(0,-1,-1)\),
\(\alpha_i=\delta_0(q_i)\), and
\(\theta=-\delta_1^{-1}(\alpha_1\smile\alpha_2)\).  The standard
\v Cech total differential gives exactly
\[
 \delta_1\bigl(\rho_*(\alpha\otimes\eta_i)\bigr)
 =-\alpha\smile\alpha_i,
\]
and therefore the surjective row \((0,1,-1,0)\) in the ordered basis
\((\alpha_1\eta_1,\alpha_1\eta_2,\alpha_2\eta_1,\alpha_2\eta_2)\).
}

\ResultHead{The unique global lift, invariant frame, and scalar}
{\ResultSmall
The lift of \(1\) through
\(0\to\mathcal O(-1)\to f_*\mathcal E\to\mathcal O\to0\) is unique because
\(H^0(\mathcal O(-1))=0\).  On the marked cover it is
\[
 h=(a(z),1)\binom{d\zeta_1}{d\zeta_2},\qquad a(z_2)=\frac12.
\]
The genuine order-four invariant target frame is
\[
 \omega_2=s\,u_2(s)(\gamma\wedge u)\ell,\qquad u_2(0)=1;
\]
the uncorrected \(s(\gamma\wedge u)\ell\) is invariant only at \(s=0\).
Consequently
\[
 \kappa_1(h\otimes\gamma)(p_2)
 =-\frac14\bigl(a(z_2)\dot\tau(z_2)+\dot\mu(z_2)\bigr)
 =-\frac14\dot\mu(z_2)\ne0,
\]
because \(\dot\tau(z_2)=0\) and
\(K(z_2)=6\dot\mu(z_2)^2\ne0\).
}

\ResultHead{Bott--Chern normalization and scope}
{\ResultSmall
Orient \(\omega=c_1(\mathcal O_B(1))\).  Then
\[
 \exp(12\mathfrak g)=f^*\mathcal O_B(1),\qquad
 c_{1,BC}(\exp(z\mathfrak g))=z\mathsf A+\bar z\,\mathsf B,
\]
so \(c_{1,BC}(f^*\mathcal O_B(1))=12(\mathsf A+\mathsf B)\).
Together with the unit \(\mathbf1=[1]_{BC}\), the nine positive-degree
generators are bases and span a square-zero ideal \(N\).  The unit statement
is literal: a \(d\)-closed function on compact connected \(X\) is constant,
and there is no \(\partial\bar\partial\) image of bidegree \((0,0)\), so
\(H^{0,0}_{BC}(X)=\mathbb C[1]\).  The exact fibre
pairings above prove the required Dolbeault nonvanishing.  The proof of
\(\mathsf A\mathsf C=0\) additionally uses Dolbeault--Serre duality and the
independently proved injection
\(H^{3,2}_{BC}(X)\hookrightarrow H^{3,2}_{\bar\partial}(X)\); the
Bott--Chern--Aeppli perfect pairing alone does not annihilate that product.
No preferred literal coefficient atlas for the noncanonical global smooth
forms is claimed.
}

\vfill
{\fontsize{6.25pt}{7.1pt}\selectfont\color{AuditProvenance}
The exact proofs are
\texttt{candidate\_geometry.tex}, \texttt{hodge\_frolicher.tex}, and
\texttt{bott\_chern\_algebra.tex}.}
\endgroup

\newpage
\thispagestyle{empty}
\begingroup
\raggedright
\setlength{\parindent}{0pt}
\setlength{\parskip}{0pt}
\setlength{\fboxsep}{5pt}
\newcommand{\ResultHead}[1]{\par\vspace{2.7pt}{\fontsize{8.4pt}{9.25pt}\selectfont\bfseries\color{AuditReviewed}#1}\par\vspace{0.8pt}}
\newcommand{\ResultSmall}{\fontsize{7.05pt}{8.05pt}\selectfont}

{\fontsize{14.6pt}{16pt}\selectfont\bfseries Typed affine-defect and toric restriction certificates}
\par\vspace{1pt}
{\fontsize{6.8pt}{7.75pt}\selectfont\color{AuditProvenance}
This page records the final type, sign, inclusion-direction, and antecedent
repairs in the purely analytic all-Picard proof.}

\ResultHead{The finite defect identity is typed at every cover level}
{\ResultSmall
Let
\[
 r_j:\widetilde J_j=\Delta_j\times\mathbb C^2\longrightarrow J'_j
\]
be the lattice quotient.  On the vector cover put
\[
 Q_{j,\log}=\widetilde g_j^{\,v_j/m_j},\qquad
 Q_{j,\mathrm{lin}}=\widetilde g_j^{\,0}.
\]
The exact clockwise relations are
\[
 Q_{j,\log}\tau_\lambda Q_{j,\log}^{-1}=\tau_{A_j\lambda},
 \quad (Q_{j,\log})^{m_j}=\tau_{v_j},\qquad
 Q_{j,\mathrm{lin}}\tau_\lambda Q_{j,\mathrm{lin}}^{-1}=\tau_{A_j\lambda},
 \quad (Q_{j,\mathrm{lin}})^{m_j}=1.
\]
The positive \(v_j\) sign follows from
\(s_j\circ g_j=e^{-2\pi i/m_j}s_j\).  The descended
\(q_{j,\log},q_{j,\mathrm{lin}}\) on \(J'_j\) have order dividing \(m_j\)
because \(\tau_{v_j}\) dies in the lattice quotient, and order exactly
\(m_j\) because both cover \(g_j\), whose order is \(m_j\).
On \(J_j'^*=J'_j|_{\Delta_j^*}\), the branch-independent translation and
the two quotient maps give the commuting square
\[
 h_jq_{j,\log}h_j^{-1}=q_{j,\mathrm{lin}},\qquad
 p_{j,\mathrm{lin}}\circ h_j=\bar h_j\circ p_{j,\log},
\]
where
\[
 p_{j,\log}:J_j'^*\to N_j^*,\qquad
 p_{j,\mathrm{lin}}:J_j'^*\to J^*,\qquad
 \bar h_j:N_j^*\xrightarrow{\ \sim\ }J^*.
\]
The line-bundle comparison is exactly
\[
 p_{j,\log}^*(L|_{N_j^*})
 \simeq h_j^*p_{j,\mathrm{lin}}^*(L|_{J^*}).
\tag{5.17-A14a}
\]
For the factor extension return to the filled vector cover:
\[
 P_{j,\log}:J'_j\to N_j,\qquad
 \pi_j^{\rm vec}:=P_{j,\log}\circ r_j:\widetilde J_j\to N_j,
\]
where \(P_{j,\log}\) is the full quotient and
\(P_{j,\log}|_{J_j^{\prime *}}=p_{j,\log}\).  Its deck group is
\[
 \Gamma_j^{\rm log}
 =\langle\tau_\lambda,Q_{j,\log}:\lambda\in\Lambda\rangle.
\]
The pullback \((\pi_j^{\rm vec})^*(L|_{N_j})\) is trivial on the
contractible Stein space \(\widetilde J_j\); its factors define an integral
central extension, pushed out along \(\mathbb Z\hookrightarrow\mathbb R\).
Over \(\Delta_j^*\), the line-bundle square \textup{(5.17-A14a)} identifies
each fibre with the global factor extension by a translation, so its
commutator form is \(E\).  This form is integer-valued and locally constant
on the connected disc \(\Delta_j\), hence remains \(E\) on the filled cover.
Thus one may choose normalized vertical lifts satisfying
\[
 U_\lambda^{(j)}U_\mu^{(j)}
 =\mathbf c\!\left(\frac12E(\lambda,\mu)\right)
  U_{\lambda+\mu}^{(j)}.
\]
For \(\alpha\in\Gamma_j^{\log}\), define \(A_\alpha,T_\alpha\) by
\[
 \alpha\tau_\lambda\alpha^{-1}=\tau_{A_\alpha\lambda},\qquad
 T_\alpha=(A_\alpha^{-1})^{\mathsf t},
\]
and define
\(\widetilde d_j(\alpha)\in V_{\mathbb R}\) by
\[
 \widetilde G_\alpha
 U_{A_\alpha^{-1}\lambda}^{(j)}\widetilde G_\alpha^{-1}
 =\mathbf c\!\left(\left\langle\widetilde d_j(\alpha),\lambda\right\rangle\right)
 U_\lambda^{(j)}.
\tag{5.17-A14c}
\]
\[
 \widetilde d_j(\alpha\beta)
 =\widetilde d_j(\alpha)+T_\alpha\widetilde d_j(\beta).
\tag{5.17-A14d}
\]
The central exponent is additive in \(\lambda\), which makes
\(\widetilde d_j(\alpha)\) unique; composing lifts gives the displayed
crossed law.  Put
\[
 \widetilde J_j^*=\Delta_j^*\times\mathbb C^2,
\]
\[
\begin{aligned}
 \Phi_{j,\log}
  &:=p_{j,\log}\circ r_j|_{\widetilde J_j^*},\\
 \Phi_{j,\mathrm{glob}}
  &:=p_{j,\mathrm{lin}}\circ h_j\circ r_j|_{\widetilde J_j^*}
   =\bar h_j\circ\Phi_{j,\log}.
\end{aligned}
\]
These maps are \(\Gamma_j^{\log}\)-invariant.  Pullback of
\textup{(5.17-A14a)} to this common cover gives
\[
 \Theta_j:
 \Phi_{j,\log}^*(L|_{N_j^*})
 \xrightarrow{\ \sim\ }
 \Phi_{j,\mathrm{glob}}^*(L|_{J^*}),
\]
and hence a centre-preserving isomorphism of the pushed-out
\(\Gamma_j^{\log}\)-factor extensions.

Let \(\eta_j\) be the total-space loop obtained by applying the zero section
to the based clockwise meridian \(\rho_j\).  Choose
\(x_j\in J_j^{\prime *}\) over \(\eta_j(0)\), and put
\(y_j=h_j^{-1}(x_j)\).  The \(p_{j,\mathrm{lin}}\)-lift of \(\eta_j\)
from \(x_j\) ends at
\(q_{j,\mathrm{lin}}(x_j)\); its \(h_j^{-1}\)-transport is the
\(p_{j,\log}\)-lift of \(\bar h_j^{-1}\eta_j\) from \(y_j\), ending at
\(q_{j,\log}(y_j)\).
For a vector lift \(\widetilde y_j\), the adjacent branch
\[
 \ell'(w)=\ell(g_j^{-1}w)-\frac{2\pi i}{m_j}
\]
has endpoint exactly \(Q_{j,\log}(\widetilde y_j)\), with no inverse or
lattice translate.  Thus, for some \(a_j\in V_{\mathbb R}\) and
\(b_j\in\mathbb R\),
\[
 \Theta_{j,*}(U_\lambda^{(j)})
 =\mathbf c(\langle a_j,\lambda\rangle)U_\lambda,\qquad
 \Theta_{j,*}(\widetilde G_{Q_{j,\log}})
 =\mathbf c(b_j)G_{\rho_j}.
\]
The central factor \(\mathbf c(b_j)\) drops out under conjugation, so
\[
 \widetilde d_j(Q_{j,\log})=d_j+(T_j-I)a_j,
 \qquad
 \widetilde d_j(\tau_v)=\iota_vE,
 \qquad (\iota_vE)(\lambda)=E(v,\lambda).
\]
Since \(N_{T_j}(T_j-I)=T_j^{m_j}-I=0\),
\[
 N_{T_j}d_j
 =\widetilde d_j((Q_{j,\log})^{m_j})
 =\widetilde d_j(\tau_{v_j})=\iota_{v_j}E.
\tag{5.17-A15}
\]
Under \(\ell\mapsto\ell+2\pi in\),
\[
 H_j\longmapsto\tau_{nv_j}H_j,\qquad
 \Theta'_{j,*}=\operatorname{Ad}(U_{nv_j})\circ\Theta_{j,*},\qquad
 a_j\longmapsto a_j+n\iota_{v_j}E.
\]
Since \(A_jv_j=v_j\) and \(E\) is \(A_j\)-invariant,
\(T_j(\iota_{v_j}E)=\iota_{v_j}E\).  Thus (5.17-A14e) and
(5.17-A15) are branch- and choice-independent and fix both type and sign.
}

\ResultHead{The toric affine opens and descent tensor product}
{\ResultSmall
For adjacent rays in a maximal cone \(\sigma=\{v,w,u\}\), the inclusions are
\[
 U_{\rho_v}=D(z_wz_u)\subset U_\sigma,\qquad
 U_{\rho_w}=D(z_vz_u)\subset U_\sigma,
 \qquad U_{\rho_v}\cap U_{\rho_w}=D(z_vz_wz_u)=U_0.
\]
Thus the maximal affine chart is not contained in the ray-chart
intersection.  Instead the contractible Stein chart
\[
 S_\sigma=\{(z_v,z_w,z_u)\in\mathbb C^3:
 |z_vz_wz_u|<\epsilon\}
\]
has \(\operatorname{Pic}(S_\sigma)=0\) and yields the required
zero-horizontal-winding frame comparison.  If
\[
 \mu_\lambda^+:\Psi_\lambda^*\mathcal O(D_a)
 \longrightarrow\mathcal O(D_a),\qquad
 s\longmapsto\chi^{m_\lambda}s,
\]
then the map required on the negative divisor by the pullback tensor product
is the inverse-dual
\[
 \mu_\lambda^-=((\mu_\lambda^+)^{-1})^\vee:
 \Psi_\lambda^*\mathcal O(-D_a)\longrightarrow\mathcal O(-D_a),
 \qquad s\longmapsto\chi^{-m_\lambda}s.
\]
The ordinary dual has the opposite arrow and is not the required descent
map.
}

\ResultHead{The local bundle, its form, and the exact antecedent}
{\ResultSmall
Put \(M=L|_{N_0}\), let \(E_M\) be the Appell--Humbert form of
\(M|_{F_t}\), and fix the typed lattice isomorphism
\[
 \Lambda_{\rm tor}:=\ker(M_0-I)=\operatorname{im}(M_0-I)
 =\mathbb Z\langle\widehat w,\widehat\delta\rangle,
 \qquad \overline\Lambda:=\Lambda/\Lambda_{\rm tor},
\]
\[
 B_0:\overline\Lambda\xrightarrow{\ \sim\ }\Lambda_{\rm tor},
 \qquad
 B_0\overline{\widehat\gamma}=-\widehat\delta,\qquad
 B_0\overline{\widehat u}=\widehat w.
\]
Its matrix is
\(\left(\begin{smallmatrix}0&1\\-1&0\end{smallmatrix}\right)\)
from \((\overline{\widehat\gamma},\overline{\widehat u})\) to
\((\widehat w,\widehat\delta)\).
If \((d_1,d_2,d_3)\) are the independent degrees of
\(\nu^*(M|_W)\), then all six degrees are
\[
 (d_1,d_2,d_3,d_1,d_2,d_3),
\]
and the support-function polarization satisfies
\[
 B=\begin{pmatrix}d_2+d_3&d_3\\d_3&d_1+d_3\end{pmatrix},
 \qquad
 E_M((\lambda,0),(0,m))=-B(B_0\lambda,m)
 \quad(\lambda\in\overline\Lambda,\ m\in\Lambda_{\rm tor}).
\]
The preceding finite-defect calculation proves \(k=0\), hence
\(E_M=c_1(M|_{F_t})=c_1(L|_{F_t})=0\).  Only with this antecedent does
unimodularity of
\(B_0\) force \(B=0\); the displayed matrix then forces
\(d_1=d_2=d_3=0\), and all six degrees vanish.  Therefore
\(\nu^*(M|_W)\simeq\mathcal O_{\mathrm{dP}_6}\), proving the global
restriction map zero.
}

\vfill
{\fontsize{6.25pt}{7.1pt}\selectfont\color{AuditProvenance}
The complete proof, including the argument-principle cusp value and residual
frame winding, is \texttt{analytic\_picard\_completion.tex}.  It uses neither
global integral Leray transgression nor local integral nearby-cycle
specialization.  Independent specialist replication remains open.}
\endgroup

\newpage
\thispagestyle{empty}
\begingroup
\raggedright
\setlength{\parindent}{0pt}
\setlength{\parskip}{0pt}
\setlength{\fboxsep}{5pt}
\newcommand{\ResultHead}[1]{\par\vspace{2.5pt}{\fontsize{8.35pt}{9.2pt}\selectfont\bfseries\color{AuditReviewed}#1}\par\vspace{0.8pt}}
\newcommand{\ResultSmall}{\fontsize{7.05pt}{8.08pt}\selectfont}

{\fontsize{14.6pt}{16pt}\selectfont\bfseries 3 September 2026: native cusp, sphere carriers, and spectral scope}
\par\vspace{1pt}
{\fontsize{6.8pt}{7.75pt}\selectfont\color{AuditProvenance}
This is a later workbench update.  It records the precise answer to the
separate question whether the native cusp degeneration can be assembled into
an \(\mathbb R^4\)-bearing construction with source \(S^{12}\) or \(S^{24}\).}

\ResultHead{The native degeneration is \(\mathbb T^2\times\mathbb R^2\) --- VERIFIED}
{\ResultSmall
On the positive radial cusp, the marked fibre is
\[
 F_T=\mathbb R^4/P_T\mathbb Z^4,\qquad
 \det P_T=-D_T,\qquad D_T=T^2+(2H-B)T+O(1).
\]
The periods \(P_T\widehat w=e_1\) and
\(P_T\widehat\delta=e_3\) have length one for every \(T\).  With
\(K_\Lambda=\mathbb Z\langle\widehat w,\widehat\delta\rangle\), projection
to the other two real coordinates gives the exact connected quotient
\[
 0\longrightarrow\mathbb T^2\longrightarrow F_T
 \xrightarrow{\ \pi_T\ }Q_T\longrightarrow0,
 \qquad Q_T=\mathbb R^2/Y_T\mathbb Z^2,
 \]
where \(Y_T=T\left(\begin{smallmatrix}0&1\\-1&0\end{smallmatrix}\right)+O(1)\).
Thus \(\operatorname{sys}(Q_T)\ge T-C\) and the unmodified fibres converge
pointed locally isometrically to
\(\mathbb T^2\times\mathbb R^2\).  This statement is intrinsic to the
candidate's original \(S^6\) cusp; no auxiliary cover or second manifold is
used.
}

\ResultHead{The finite-cover \(\mathbb R^4\) limit is valid but is a different construction}
{\ResultSmall
After \(t=v^N\), the corrected source multiplier is
\(C_N(v)=C(v^N)/N\).  Toric multiplication then descends across the central
fibre to
\[
 \kappa_N:N_0^{\langle N\rangle}\longrightarrow N_0,\qquad
 \deg\kappa_N=N^3,\qquad \kappa_N^*W=NW_N.
\]
Factoring through the normal base change gives a generic degree-\(N^2\)
map, quasi-etale but nonflat precisely over the double and triple strata.
Its regular-fibre lattice replaces the two unit periods by
\(Ne_1,Ne_3\), so \(T,N\to\infty\) gives a pointed local
\(\mathbb R^4\) limit.  This is an auxiliary local Kummer cover of the cusp,
not a global \(S^{12}/S^{24}\) construction and not a cover of \(X\).
}

\ResultHead{The exact \(S^{12}\) and \(S^{24}\) answer uses smash sources and conical targets}
{\ResultSmall
The native quotient exists on the whole punctured cusp because \(M_0\) is
the identity on \(\Lambda/K_\Lambda\).  Collapsing the central fibre and the
exterior of a closed cusp neighborhood gives a based quotient
\[
 q_2:X\longrightarrow\mathcal A=\operatorname{Susp}(T^3),\qquad
 q_2^{\wedge2}:X^{\wedge2}\cong S^{12}
 \longrightarrow\mathcal A^{\wedge2}.
\]
On two synchronized radial fibres its restriction is exactly
\(\pi_T\times\pi_T:F_T^2\to Q_T^2\), and the intrinsic flat tori
\(Q_T^2\) converge pointed locally isometrically to \(\mathbb R^4\).

The unique primitive character invariant under all three monodromies extends
over the order-three and order-four fillings to explicit Seifert targets.
After coning the remaining cusp boundary torus one obtains
\[
 q_1:X\longrightarrow\mathcal C,\qquad
 q_1^{\wedge4}:X^{\wedge4}\cong S^{24}
 \longrightarrow\mathcal C^{\wedge4}.
\]
Its synchronized radial restriction is four expanding circles
\(C_T^4\), with
\(\operatorname{sys}(C_T^4)=T+H-B+O(T^{-1})\), hence again an intrinsic
pointed \(\mathbb R^4\) limit.  These are genuine global maps with real
sphere sources.  Their targets have nonspherical integral homology and are
conically stratified; in every fixed compatible ambient target metric the
radial slices collapse to the cone vertex.  They therefore construct neither
a complex structure on \(S^{12}\) or \(S^{24}\) nor a fixed compact
four-dimensional spacetime.
}

\ResultHead{What the finite gauge calculation proves --- and what remains OPEN}
{\ResultSmall
With independent cover index \(M\) and cubical mesh \(N\), an explicit
periodic \(\mathrm{SU}(2)\) link field has
\[
 U_{12}(n)=\exp(2\pi H/N^2),\qquad
 S^{\log}_{T,M,N}=\frac{2\pi^2M^2}{g_{\rm YM}^2D_T}.
\]
The exactly diagonalized finite Coulomb oscillator has gap
\[
 \min_{k\ne0}\frac1{g_{\rm YM}}
 \bigl\lVert P_T^{-\mathsf T}S_M^{-1}q_N(k)\bigr\rVert.
\]
The sequence \(T=j^2,M=j,N=j^3\) simultaneously gives local
\(\mathbb R^4\) cover geometry, mesh diameter \(O(j^{-1})\), Gaussian gap
\(O(j^{-2})\), and logarithmic action \(O(j^{-2})\).  This is a theorem for
a finite quadratic oscillator on a linear Coulomb space.  No equality with
the compact-link Haar/Gauss-law Hilbert space, no transfer-matrix or
Osterwalder--Schrader reconstruction, and no interacting four-dimensional
Yang--Mills mass-gap conclusion has been proved.
}

\vfill
{\fontsize{6.25pt}{7.1pt}\selectfont\color{AuditProvenance}
Complete proofs are in the four cusp sections indexed by
\nolinkurl{supporting_materials/workbench/research/README.md}.  The preceding
eight pages retain the earlier review and validation record unchanged.}
\endgroup
